\documentclass[twocolumn,prl,aps,floatfix]{revtex4-2}
\usepackage{amsmath,amssymb,graphicx,bm,xcolor}
\usepackage[colorlinks=true,linkcolor=blue,citecolor=blue,urlcolor=blue]{hyperref}
\graphicspath{{../figures/}}

\begin{document}

\title{GRAPE-Like Optimization of Dispersive Readout Pulses in Circuit~QED:\\
A Systematic Study of Objectives, Waveform Families, and Robustness}

\author{Readout Pulse Optimization Study}
\affiliation{cQED Autonomous Research Platform}

\date{\today}

\begin{abstract}
We perform a systematic study of readout pulse optimization in circuit quantum
electrodynamics (cQED) using the linear dispersive model for the qubit-state-conditioned
resonator field.  We compare five pulse families (square, Gaussian, Hann window,
cosine-rise, and spline) against GRAPE-like piecewise-constant control across a sweep of
dispersive coupling ratios $\chi/\kappa \in \{0.25, 0.5, 1.0, 2.0, 4.0\}$ and
dimensionless readout durations $\kappa T \in [0.5, 20]$.  We find that: (i) driving
at the midpoint $\omega_d = \omega_r + \chi/2$ is optimal for all waveform families;
(ii) the time-integrated SNR$^2$ is the appropriate objective, not endpoint pointer
separation; (iii) after enforcing a common peak-amplitude bound, the square pulse is the
best analytical family across the tested parameter range; (iv) GRAPE converges to that
same square-like solution within numerical precision and does not provide a meaningful
SNR$^2$ advantage in this linear, peak-limited setting; (v) the current single-segment
active-depletion heuristic does not outperform passive ring-down; (vi) the optimal readout
duration balancing SNR gain against $T_1$-induced infidelity occurs at
$\kappa T \approx 5.79$ ($T \approx 184\,$ns for $\kappa/2\pi = 5\,$MHz), where
$F_{\rm net}^{\rm max} = 0.9967$; and (vii) the current robust-GRAPE formulation sacrifices
nominal performance without improving on the square-pulse baseline. We therefore recommend
a midpoint-driven square pulse as the practical default for this model class, with GRAPE
reserved for future formulations that include additional hardware or physics constraints.
\end{abstract}

\maketitle

% ─────────────────────────────────────────────────────────────────────────────
\section{Introduction}
% ─────────────────────────────────────────────────────────────────────────────

Dispersive readout is the canonical method for measuring superconducting qubits in the
circuit QED architecture~\cite{blais2021circuit,krantz2019quantum}.  A microwave tone
drives a readout resonator coupled dispersively to the qubit; the qubit-state-dependent
resonator response encodes the qubit state in the phase and amplitude of the transmitted
signal.  While the physics of dispersive readout is well understood
\cite{gambetta2007protocols,boissonneault2010dispersive}, the optimal control
of the readout pulse shape has received less systematic study than gate pulse design.

Three practical requirements pull in different directions:
(i) \emph{fast} readout (large signal-to-noise ratio per unit time),
(ii) \emph{clean} reset (minimal residual resonator photons after readout), and
(iii) \emph{QND} character (minimal qubit disturbance during readout).
Optimizing these simultaneously calls for a carefully chosen objective function,
waveform class, and control formulation.

GRAPE-like optimal control~\cite{khaneja2005optimal} has been applied widely to gate
pulse design in cQED but is less commonly used for readout.  Active resonator
reset~\cite{reed2010fast,McClure2016rapid} and fast readout by waveform shaping
\cite{walter2017rapid,jeffrey2014fast} have been demonstrated experimentally, but a
systematic comparison of the problem formulations in the readout context is lacking.

This study fills that gap. We deliberately avoid assuming that GRAPE is the right answer
from the outset, and instead treat the choice of objective function, waveform family,
and optimization method as the central research questions.

% ─────────────────────────────────────────────────────────────────────────────
\section{Physical Model}
% ─────────────────────────────────────────────────────────────────────────────

\subsection{Linear dispersive resonator model}

In the frame rotating at the drive frequency $\omega_d$, the qubit-state-conditioned
resonator coherent amplitude $\alpha_j(t)$ ($j \in \{g, e\}$) obeys
\begin{equation}
  \dot{\alpha}_j = -\!\left(\frac{\kappa}{2} + i\Delta_j\right)\alpha_j - i\epsilon(t),
  \label{eq:ode}
\end{equation}
where $\kappa$ is the resonator decay rate, $\Delta_g = \omega_r - \omega_d$ is the
ground-state detuning, $\Delta_e = \omega_r + \chi - \omega_d$ is the excited-state
detuning, $\chi$ is the dispersive shift, and $\epsilon(t)$ is the drive envelope
(complex, units rad/s).  The output field in input-output theory is
$s_j(t) = \sqrt{\kappa}\,\alpha_j(t)$.

\subsection{Steady-state analysis}

At steady state ($\dot{\alpha} = 0$):
\begin{equation}
  \alpha_j^{\rm ss} = \frac{-i\epsilon}{\kappa/2 + i\Delta_j}.
  \label{eq:ss}
\end{equation}
The steady-state pointer separation is maximized at the \emph{midpoint drive}
$\omega_d = \omega_r + \chi/2$ ($\Delta_g = -\chi/2$, $\Delta_e = +\chi/2$):
\begin{equation}
  |\Delta\alpha^{\rm ss}_{\rm max}| = \frac{4|\epsilon||\chi|}{\kappa^2 + \chi^2},
  \label{eq:sep_formula}
\end{equation}
which peaks at $|\chi/\kappa| = 1$ with value $2|\epsilon|/\kappa$.  This identifies
$|\chi| = \kappa$ as the optimal dispersive coupling ratio.

\subsection{Physical parameters}

We study parameters representative of state-of-the-art transmon experiments:
$\kappa/2\pi = 5\,\mathrm{MHz}$, $|\chi|/2\pi \in \{1.25, 2.5, 5.0, 10.0, 20.0\}\,\mathrm{MHz}$
(corresponding to $|\chi|/\kappa \in \{0.25, 0.5, 1.0, 2.0, 4.0\}$),
$|\epsilon_{\rm max}|/2\pi = 5\,\mathrm{MHz}$ (corresponding to $\lesssim 5$ steady-state
photons on resonance), $T_1 = 30\,\mu\mathrm{s}$, and $T_2 = 20\,\mu\mathrm{s}$.

% ─────────────────────────────────────────────────────────────────────────────
\section{Objective Functions}
% ─────────────────────────────────────────────────────────────────────────────

\subsection{Objective A: Endpoint pointer separation}
\begin{equation}
  J_A = -|\alpha_e(T) - \alpha_g(T)|^2.
\end{equation}
This captures the \emph{final} phase-space separation but ignores the transient
signal accumulated during ring-up.  Our numerical study confirms that $J_A$ saturates
at $|\Delta\alpha^{\rm ss}|^2$ well before $J_B$ saturates, making it a poor proxy
for SNR at intermediate $\kappa T$.

\subsection{Objective B: Time-integrated distinguishability (SNR$^2$)}
\begin{equation}
  J_B = -\kappa \int_0^T |\alpha_e(t) - \alpha_g(t)|^2 \,dt = -{\rm SNR}^2.
  \label{eq:jb}
\end{equation}
This equals $-\int_0^T |s_e(t) - s_g(t)|^2\,dt$, the integrated output-field
distinguishability.  For quantum-limited heterodyne detection, SNR$^2$ directly
determines the assignment fidelity:
\begin{equation}
  F_A = 1 - \tfrac{1}{2}\,{\rm erfc}({\rm SNR}/\sqrt{2}),
  \label{eq:fa}
\end{equation}
making $J_B$ the most physically grounded single-objective function.

\subsection{Objective C: Multi-objective with residual photons}
\begin{equation}
  \mathcal{L} = -w_1\,{\rm SNR}^2 + w_2\,n_{\rm res}(T),
  \label{eq:multi}
\end{equation}
where $n_{\rm res}(T) = \tfrac{1}{2}(|\alpha_e(T)|^2 + |\alpha_g(T)|^2)$ is the
mean residual photon number.  We study the Pareto tradeoff between SNR$^2$ and
$n_{\rm res}$ by sweeping $w_2/w_1$.

% ─────────────────────────────────────────────────────────────────────────────
\section{GRAPE for the Linear Dispersive Model}
% ─────────────────────────────────────────────────────────────────────────────

For a piecewise-constant drive $\epsilon(t) = \epsilon_k$ for $t \in [k\tau, (k+1)\tau)$
with $\tau = T/N$, the propagator per segment is
\begin{align}
  \alpha_j(t_{k+1}) &= A_j\,\alpha_j(t_k) + B_j\,\epsilon_k, \\
  A_j &= e^{-\lambda_j\tau},\quad
  B_j = \frac{-i(A_j - 1)}{\lambda_j},\quad
  \lambda_j = \tfrac{\kappa}{2} + i\Delta_j.
\end{align}
The Wirtinger gradient of ${\rm SNR}^2$ with respect to $\epsilon_k^*$ is computed
by the adjoint (backward) recursion:
\begin{equation}
  \frac{\partial\,{\rm SNR}^2}{\partial \epsilon_k^*}
    = \kappa\tau\bigl(B_e\,\Lambda_e^{(k)} - B_g\,\Lambda_g^{(k)}\bigr),
  \label{eq:grape_grad}
\end{equation}
where $\Lambda_j^{(k)} = \bigl(\Delta\alpha(t_{k+1})\bigr)^* + A_j\,\Lambda_j^{(k+1)}$
with $\Lambda_j^{(N)} = 0$ (backward from $k = N-1$ to $0$).
This is the standard GRAPE adjoint method applied to the linear dispersive ODE.

Optimization is performed with L-BFGS-B using $N = 60$ segments, 10 random restarts,
and amplitude bound $|\epsilon_k| \leq |\epsilon_{\rm max}|$.

% ─────────────────────────────────────────────────────────────────────────────
\section{Results}
% ─────────────────────────────────────────────────────────────────────────────

\subsection{Phase 1: Analytical theory}

Figure~\ref{fig:steady_state} shows $|\Delta\alpha^{\rm ss}_{\rm max}|$ vs $|\chi|/\kappa$
(panel a) and vs drive detuning $\Delta_g$ (panel b).  The analytical formula
\eqref{eq:sep_formula} is confirmed numerically to better than $10^{-6}$ relative error,
and agrees with \texttt{cqed\_sim.ReadoutResonator} to machine precision.  The midpoint
drive $\Delta_g = -\chi/2$ is confirmed as optimal.

\begin{figure}[t]
  \centering
  \includegraphics[width=\columnwidth]{fig1_steady_state.pdf}
  \caption{%
    (a)~Steady-state pointer separation $|\Delta\alpha^{\rm ss}_{\rm max}|$ vs $|\chi|/\kappa$.
    The maximum occurs at $|\chi|/\kappa = 1$ where $|\Delta\alpha^{\rm ss}| = 2|\epsilon|/\kappa$.
    (b)~$|\Delta\alpha^{\rm ss}|$ vs drive detuning $\Delta_g$ at $|\chi|/\kappa = 1$,
    confirming the midpoint optimum at $\Delta_g = -\chi/2$.
  }
  \label{fig:steady_state}
\end{figure}

Figure~\ref{fig:snr_vs_kT} shows SNR$^2$ and assignment fidelity $F_A$ vs $\kappa T$
for midpoint and on-resonance drives.  The midpoint drive outperforms the on-resonance
drive by $50$--$100\%$ at intermediate $\kappa T$, converging at large $\kappa T$
as both approach steady state.

\begin{figure}[t]
  \centering
  \includegraphics[width=\columnwidth]{fig2_snr_vs_kappa_t.pdf}
  \caption{%
    SNR$^2$ (a) and assignment fidelity $F_A$ (b) vs $\kappa T$ for square pulses at the
    midpoint drive ($\Delta_g = -\chi/2$, blue) and on-resonance ($\Delta_g = 0$, red).
    The steady-state approximation $\kappa T|\Delta\alpha^{\rm ss}|^2$ (dotted) overestimates
    SNR$^2$ at small $\kappa T$ where the resonator has not yet rung up.
  }
  \label{fig:snr_vs_kT}
\end{figure}

\subsection{Phase 2: Pulse-family baseline}

Figure~\ref{fig:families} shows SNR$^2$ vs $|\chi|/\kappa$ at $\kappa T = 5$ (panel a)
and vs $\kappa T$ at $|\chi|/\kappa = 1$ (panel b) for all five pulse families under a
shared \emph{peak}-amplitude bound.  Under that fair comparison, the square pulse is the
best-performing analytical family throughout the tested $|\chi|/\kappa$ range.  The
cosine-rise pulse remains competitive but never exceeds square, while Gaussian, Hann, and
spline envelopes underperform because their smooth turn-on uses less of the available peak
drive budget in the strictly linear model.

The key finding is that \textbf{the midpoint-driven square pulse saturates the best
simple-family performance} in this study.  At $|\chi|/\kappa = 1$ and $\kappa T = 5$,
the square pulse yields SNR$^2 = 8.93$, whereas the Hann envelope yields only
SNR$^2 = 3.19$.

Active depletion (Fig.~\ref{fig:depletion}) should be treated as a negative result here:
the implemented single-segment heuristic does not beat passive ring-down and therefore is
not yet a useful control primitive.

\begin{figure}[t]
  \centering
  \includegraphics[width=\columnwidth]{fig5_depletion.pdf}
  \caption{%
    Residual photon number after readout as a function of additional depletion time
    $\kappa T_{\rm dep}$.  The dotted line shows the readout-only residual occupancy,
    the red curve shows passive natural ring-down, and the blue curve shows the
    implemented active-depletion heuristic.  In this study the heuristic fails to beat
    passive ring-down and is therefore not recommended.
  }
  \label{fig:depletion}
\end{figure}

\begin{figure}[t]
  \centering
  \includegraphics[width=\columnwidth]{fig4_pulse_families.pdf}
  \caption{%
    SNR$^2$ vs $|\chi|/\kappa$ at $\kappa T = 5$ (a) and vs $\kappa T$ at $|\chi|/\kappa = 1$ (b)
    for five pulse families under a shared peak-amplitude bound.  The square pulse is the
    best analytical family across the tested regime; smoother envelopes sacrifice SNR$^2$
    because they do not fully use the allowed peak drive.
  }
  \label{fig:families}
\end{figure}

\subsection{Phase 3: GRAPE optimization}

Figure~\ref{fig:grape}(a) compares GRAPE against all pulse families across $|\chi|/\kappa$.
Once the waveform-family comparison is corrected to obey the same peak-amplitude limit,
GRAPE no longer outperforms the square baseline in any meaningful way: the optimized
SNR$^2$ values match the square pulse to within numerical precision across the full
$|\chi|/\kappa$ sweep.

Figure~\ref{fig:grape}(b) shows the same pattern across $\kappa T$: GRAPE and square are
indistinguishable within numerical precision from $\kappa T = 1$ to $12$, while the Hann
envelope remains substantially worse under the same bound.  Segment convergence remains
excellent, with SNR$^2$ stable to better than $10^{-4}$ between $N = 60$ and $N = 120$
segments, confirming that the optimizer itself is numerically well behaved.

The multi-objective Pareto sweep provides another negative result: sweeping the residual-
photon penalty over $w_2/w_1 \in [0, 5]$ leaves both SNR$^2$ and $n_{\rm res}$ unchanged
within reported precision.  In the present linear model and control parameterization, the
chosen penalty does not move the optimum away from the square-like solution.

\begin{figure}[t]
  \centering
  \includegraphics[width=\columnwidth]{fig6_grape_comparison.pdf}
  \caption{%
    GRAPE vs pulse-family baselines.
    (a)~SNR$^2$ vs $|\chi|/\kappa$ at $\kappa T = 5$.
    (b)~SNR$^2$ vs $\kappa T$ at $|\chi|/\kappa = 1$.
    Under a common peak-amplitude bound, GRAPE collapses onto the square-pulse baseline
    within numerical precision across the tested regime.
  }
  \label{fig:grape}
\end{figure}

\begin{figure}[t]
  \centering
  \includegraphics[width=\columnwidth]{fig9_waveforms.pdf}
  \caption{%
    GRAPE waveform at $|\chi|/\kappa = 1$, $\kappa T = 5$:
    (a)~I/Q quadratures vs time.
    (b)~Phase-space trajectories $\alpha_{g,e}(t)$ for GRAPE and square pulse.
    (c)~Pointer separation $|\Delta\alpha(t)|$ vs time.
    (d)~Pareto frontier: SNR$^2$ vs $n_{\rm res}$ under multi-objective optimization.
    The optimized waveform is effectively square-like, and the Pareto sweep shows no
    meaningful movement away from the nominal optimum.
  }
  \label{fig:waveforms}
\end{figure}

\subsection{Phase 4: Open-system effects}

Figure~\ref{fig:open} shows the net assignment fidelity $F_{\rm net}$ and
measurement-induced dephasing $\eta_\varphi$ vs $\kappa T$.

The net fidelity is computed as
\begin{equation}
  F_{\rm net} = (1-p_{\rm decay})\,F_A({\rm SNR}^2_{\rm eff}) + p_{\rm decay}\cdot\tfrac{1}{2},
\end{equation}
where $p_{\rm decay} = 1 - e^{-T/T_1}$ and
${\rm SNR}^2_{\rm eff} = \kappa\!\int_0^T |\Delta\alpha|^2 e^{-t/T_1}\,dt$
accounts for the weight reduction from $T_1$ decay.

For our parameters ($\kappa T_1 = \kappa \times 30\,\mu{\rm s} \approx 940$), the $T_1$
correction is negligible at $\kappa T \leq 20$.  The optimal net fidelity reaches
$F_{\rm net}^{\rm max} = 0.9967$ at $\kappa T = 5.79$, corresponding to
$T = 184\,$ns for the nominal device.  The dominant limitation on useful readout duration
is therefore SNR saturation rather than $T_1$ decay.

The measurement-induced dephasing factor
$\eta_\varphi = \exp(-\int (\kappa/2)|\Delta\alpha|^2 dt)$ becomes very small during a
successful measurement; at $\kappa T = 5$ we obtain $\eta_\varphi = 0.0106$.  This is
consistent with strong pointer-state distinguishability and highlights the QND tradeoff,
not a failure of the model.

\begin{figure}[t]
  \centering
  \includegraphics[width=\columnwidth]{fig7_open_system.pdf}
  \caption{%
    (a)~Net fidelity $F_{\rm net}$ (blue, left axis) and measurement-induced dephasing
    $\eta_\varphi$ (red, right axis) vs $\kappa T$.
    (b)~SNR$^2$ and $\eta_\varphi$ vs drive amplitude $|\epsilon|/\kappa$.
    The optimal readout duration is set by SNR$^2$ saturation, not $T_1$ decay.
  }
  \label{fig:open}
\end{figure}

\subsection{Phase 5: Robustness}

Figure~\ref{fig:robust} shows 2D sensitivity maps of SNR$^2$ over $\pm 20\%$
uncertainty in $\chi$ and $\kappa$.

The square pulse and nominal GRAPE are identical within numerical precision in both nominal
and perturbed performance because the optimizer lands on the same square-like solution.
At the reference point $|\chi|/\kappa = 1$, $\kappa T = 5$, both have nominal
SNR$^2 \approx 8.93$, mean perturbed SNR$^2 \approx 8.82$, and CV $\approx 0.172$.

The current robust-GRAPE formulation performs poorly: the nominal SNR$^2$ collapses to
$0.79$ and the mean perturbed SNR$^2$ to $0.78$, with CV $0.179$, which is slightly worse
than the square baseline.  The negative result is still useful because it shows that a
naive ensemble objective can destroy the useful solution in this parameterization.

\begin{figure}[t]
  \centering
  \includegraphics[width=\columnwidth]{fig8_robustness.pdf}
  \caption{%
    SNR$^2$ sensitivity maps over $\pm 20\%$ uncertainty in $(\chi, \kappa)$ for
    four pulse designs.  Color scale: green = high SNR$^2$, red = low SNR$^2$.
    Nominal GRAPE and the square pulse are effectively identical, while the current
    robust-GRAPE formulation sacrifices most nominal performance.
  }
  \label{fig:robust}
\end{figure}

% ─────────────────────────────────────────────────────────────────────────────
\section{Discussion}
% ─────────────────────────────────────────────────────────────────────────────

\subsection{Answers to the central research questions}

\textbf{Is endpoint pointer separation too naive?}
Yes. At $\kappa T < 3$, endpoint separation significantly underestimates integrated SNR$^2$
because it ignores the transient ring-up contribution.  The integrated SNR$^2$ (Objective B)
is the correct target.

\textbf{Is the midpoint drive always optimal?}
Yes, for all waveform families studied.  The analytical result $\Delta_g = -\chi/2$
is confirmed numerically and is a robust choice.

\textbf{Are simple pulse families sufficient?}
Yes.  Under a strict peak-amplitude constraint, the square pulse is the best analytical
family in this linear model, and GRAPE matches it within numerical precision rather than
surpassing it.

\textbf{Should depletion be jointly optimized?}
Not with the present heuristic.  The implemented single-segment depletion strategy is worse
than passive ring-down, and the current multi-objective formulation does not move the
solution off the square-like optimum.

\textbf{When does GRAPE clearly add value?}
It does not add value in the present linear, peak-limited formulation.  To justify GRAPE,
future work would need additional constraints or physics such as bandwidth limits,
nonlinear resonator effects, or a richer multi-stage reset objective.

\textbf{Should the optimized pulse target nominal or robust performance?}
At present, the square pulse is the reliable choice.  The current robust-GRAPE formulation
reduces nominal performance too severely to recommend.

\textbf{Minimum realistic model complexity?}
The linear dispersive model (Level 1) is sufficient for pulse design and optimization.
The key open-system effect ($T_1$) is negligible at $\kappa T \leq 10$ for typical parameters.
The nonlinear/multi-level effects (Kerr, $\chi'$, higher transmon levels) are significant
only at very high drive amplitudes ($n_{\rm ss} \gtrsim n_{\rm crit}$) and can be checked
a posteriori.

\subsection{Limitations}

This study uses the linear dispersive model throughout.  A more complete treatment would
include: (i) Kerr nonlinearity $K_r a^{\dagger 2} a^2$ causing squeezing and deformation
of the resonator state; (ii) measurement-induced transitions
$|e\rangle \leftrightarrow |g\rangle$ (dressed-dephasing)~\cite{boissonneault2010dispersive}
at high drive powers; (iii) realistic AWG bandwidth and imperfections limiting pulse
expressivity.  These are left for future work.

% ─────────────────────────────────────────────────────────────────────────────
\section{Conclusions and Recommendations}
% ─────────────────────────────────────────────────────────────────────────────

We recommend the following implementation strategy:

\begin{enumerate}
  \item \textbf{Objective}: Use the time-integrated SNR$^2$ (Objective B).  Endpoint
    separation is an inferior proxy.
  \item \textbf{Drive frequency}: Always use the midpoint drive $\omega_d = \omega_r + \chi/2$.
    This is the single highest-value calibration to get right.
  \item \textbf{Simple baseline}: Use the square envelope at the midpoint drive.  Under the
    validated peak-amplitude constraint it is the best analytical family in this study.
  \item \textbf{GRAPE}: Do not expect a meaningful gain over square in the present linear,
    peak-limited formulation.  Reserve GRAPE for future models with additional hardware or
    physics constraints.
  \item \textbf{Joint depletion}: Do not use the current single-segment depletion heuristic as a
    production strategy.  It underperforms passive ring-down and needs redesign.
  \item \textbf{Robustness}: Prefer the square pulse over the current robust-GRAPE formulation.
    The ensemble objective used here destroys nominal performance without improving the baseline.
  \item \textbf{Open-system}: $T_1$ effects are negligible for $\kappa T \leq 10$; the practical
    operating point is $\kappa T \approx 5$--$6$ ($T \approx 184$ ns for the nominal device),
    where $F_{\rm net}$ is already near-optimal.
\end{enumerate}

\begin{acknowledgments}
This study was performed using the \texttt{cqed\_sim} framework and the autonomous research
platform developed for the Shankar Quantum Circuits Group.
\end{acknowledgments}

\bibliography{references}

\end{document}
